\chapter{TFT o TFG de desarrollo de software}
Para trabajos de desarrollo de software, es esencial estructurar el informe de manera que refleje el proceso de desarrollo completo, desde la concepción hasta la implementación y pruebas. A continuación, se detallan los puntos propuestos: 
\begin{enumerate}
    \item Introducción: Presenta el proyecto, incluyendo el contexto y la justificación del desarrollo de software propuesto. Se debe explicar la importancia del software a desarrollar y cómo se inserta en el ámbito tecnológico actual.
    \item Descripción del problema: Define el problema que el software busca resolver, detallando las necesidades o carencias que han motivado el desarrollo del proyecto. Es importante que esta sección sea clara y específica para establecer una base sólida para los objetivos y el desarrollo posterior.
    \item Objetivos: Enumera los objetivos generales y específicos del proyecto. Los objetivos deben ser concisos y medibles, proporcionando una guía clara para las etapas de desarrollo que siguen. 
    \item Marco teórico/referencial: Incluye el marco teórico y referencial que sustenta el desarrollo del software. Debe revisar conceptos, teorías, y trabajos previos relacionados con el problema a resolver, las tecnologías utilizadas, y los métodos existentes. 
    \item Metodología y plan de trabajo: Describe la metodología de desarrollo adoptada para el proyecto, incluyendo las fases de desarrollo, las técnicas y herramientas de gestión de proyectos, y el cronograma estimado. Esto establece un camino claro para el desarrollo del proyecto. 
    \item Propuesta de solución: Presenta una descripción general de la solución propuesta, incluyendo el modelo de proceso de desarrollo (por ejemplo, ágil, cascada), la arquitectura física y lógica del sistema, y las herramientas de desarrollo seleccionadas. Esta sección debe dar una visión clara de cómo se estructurará y desarrollará el software.
    \item Requerimientos: Detalla la captura y especificación de requerimientos del software, incluyendo requerimientos funcionales, no funcionales, y restricciones. Esta fase es crucial para asegurar que el software cumpla con las expectativas y necesidades de los usuarios finales.
    \item Diseño de solución: Describe el diseño del sistema, incluyendo la estructura de módulos y componentes, el esquema de la base de datos, y los prototipos de interfaz de usuario. Esta sección debe mostrar cómo se organiza el software para cumplir con los requerimientos especificados. 
    \item Implementación: Detalla cómo los componentes y módulos del sistema fueron codificados y ensamblados, incluyendo las tecnologías y patrones de diseño utilizados. Se enfoca en la traducción del diseño en software funcional, explicando la lógica de negocio, manejo de datos, y la construcción de interfaces.
    \item Pruebas del software: Expone el plan y los resultados de las pruebas realizadas, incluyendo pruebas unitarias, de integración, de sistema, y de aceptación por usuarios. Esta sección es fundamental para demostrar la calidad y funcionalidad del software.
    \item Conclusiones: Resume las contribuciones del proyecto, las limitaciones encontradas, y sugiere líneas para trabajos futuros. Esta sección reflexiona sobre el impacto del software desarrollado, tanto a nivel técnico como en la solución del problema original. 
\end{enumerate}